\documentclass[]{article}

\usepackage{amsmath}
\usepackage{multirow}
\usepackage{natbib}
\usepackage{graphicx} 


\begin{document}

Our analysis systematically compares the statistical properties of tracer transport in a chaotic point-vortex system with a canonical Lévy walk model. We begin with a qualitative overview of the trajectory morphologies and then proceed to a quantitative analysis of transport scaling, its underlying mechanisms, and the ergodicity of the ensembles.

\subsection{Tracer trajectory morphology}

A visual inspection of representative tracer trajectories reveals a strong dependence of the transport characteristics on the number of vortices, $N$, as shown in Figure \ref{fig:trajectories}. For a low vortex count ($N=5$), the tracer path is compact and exhibits quasi-periodic, looping structures. This indicates that the dynamics are weakly chaotic, and the tracer remains confined within the influence of a few dominant vortices, failing to explore the full spatial domain within the simulation time.

As $N$ increases, the trajectories become progressively more erratic and space-filling. The $N=20$ case shows intermittent behavior, with periods of local trapping punctuated by long-range excursions. This morphology intensifies in the $N=40$ system, where the trajectory consists of numerous long, nearly ballistic flights connecting distinct, tortuous regions. This visual evidence suggests a transition towards a Lévy-like transport regime at high vortex densities.

A comparison with the benchmark Lévy walk trajectories (bottom row of Figure \ref{fig:trajectories}) reinforces this observation. The strongly superdiffusive Lévy walk ($\beta=1.5$) produces a trajectory that is morphologically similar to the $N=40$ vortex case. Conversely, the normal-diffusive Lévy walk ($\beta=2.5$) results in a compact, tortuous path reminiscent of the weakly chaotic $N=5$ vortex system. A key distinction remains: the "flights" in the vortex system are continuously curved due to the persistent influence of the velocity field, unlike the perfectly straight flights of the idealized Lévy walk.

\begin{figure}[t]
    \centering
    \includegraphics[width=\columnwidth]{../input_files/plots/fig1_trajectories.png}
    \caption{Representative 2D trajectories for passive tracers in point-vortex fields (top row) and for synthetic Lévy walks (bottom row). The spatial complexity of the vortex-driven trajectories increases with the number of vortices, N, transitioning from a compact path at N=5 to a space-filling trajectory with long, curved excursions at N=40. This progression illustrates a shift towards superdiffusion, where the N=40 trajectory is morphologically similar to the superdiffusive Lévy walk with $\beta$ = 1.5, while the N=5 trajectory resembles the compact, normal-diffusive case ($\beta$ = 2.5). Trajectories are colored by normalized time from early (blue) to late (yellow).}
    \label{fig:trajectories}
\end{figure}

\subsection{Superdiffusive scaling of transport}

To quantify the transport, we calculated the ensemble-averaged mean squared displacement (MSD), $\langle \Delta r^2(\Delta t) \rangle$. Figure \ref{fig:msd} shows the MSD curves for both systems on a log-log scale. We determined the anomalous diffusion exponent, $\alpha$, by fitting a power law, $\langle \Delta r^2 \rangle \propto (\Delta t)^\alpha$, to the data.

For the point-vortex system, the results summarized in Table \ref{tab:msd_exponents} show a clear, monotonic increase in $\alpha$ with $N$. The transport transitions from near-normal diffusion for $N=5$ ($\alpha \approx 1.05$) to pronounced superdiffusion for $N=40$ ($\alpha \approx 1.52$). This confirms that increasing the vortex density enhances the chaoticity of the flow, leading to stronger anomalous transport. At short lag times ($\Delta t \lesssim 0.5$ s), all vortex systems exhibit a ballistic regime ($\alpha \approx 2$) before transitioning to their characteristic superdiffusive scaling.

The Lévy walk MSDs serve as a benchmark. The empirical exponents are in excellent agreement with the theoretical prediction $\alpha = 3 - \beta$, validating our fitting procedure. The scaling of the $N=40$ vortex system ($\alpha \approx 1.52$) is kinematically analogous to that of a Lévy walk with $\beta \approx 1.5$.

\begin{figure}[t]
    \centering
    \includegraphics[width=\columnwidth]{../input_files/plots/fig2_msd_scaling_1777306297.png}
    \caption{Ensemble-averaged Mean Squared Displacement (MSD) as a function of lag time on log-log axes for point-vortex systems (left) and Lévy walks (right). For the point-vortex model, the anomalous diffusion exponent $\alpha$, given by the slope, increases monotonically with the number of vortices N, demonstrating a transition from near-normal diffusion for N=5 to strong superdiffusion for N=40. The Lévy walk models serve as a ground truth, with their empirical MSD scaling showing excellent agreement with the theoretical prediction $\alpha = 3 - \beta$, which validates the analysis methodology.}
    \label{fig:msd}
\end{figure}

\begin{table}[h]
\centering
\caption{Anomalous diffusion exponents for the point-vortex system.}
\label{tab:msd_exponents}
\begin{tabular}{ccc}
\hline
\hline
Number of Vortices ($N$) & Exponent ($\alpha$) & $R^2$ \\
\hline
5  & $1.05$ & $0.97$ \\
10 & $1.18$ & $0.98$ \\
20 & $1.35$ & $0.97$ \\
40 & $1.52$ & $0.96$ \\
\hline
\end{tabular}
\end{table}

\subsection{Velocity correlations and memory}

The normalized velocity autocorrelation function (VACF), shown in Figure \ref{fig:vacf}, reveals the temporal memory of the velocity field experienced by the tracers. In the point-vortex system, the velocity decorrelation time, marked by the first zero-crossing of the VACF ($\tau_0$), decreases systematically as $N$ increases. The zero-crossing times are approximately $\tau_0 \approx 3.5$ s for $N=5$, $2.1$ s for $N=10$, $1.4$ s for $N=20$, and $0.9$ s for $N=40$. This reflects the faster, more chaotic evolution of the denser vortex field. Notably, the VACF for $N=20$ and $N=40$ exhibits a negative lobe, indicating velocity anti-correlations. This is a signature of the underlying Hamiltonian dynamics, where a tracer can be deflected by a vortex and sent back in a direction opposite to its previous motion.

In contrast, the VACFs for the superdiffusive Lévy walks ($\beta < 2$) show a slow, positive decay without a negative lobe. The long persistence of positive correlations is a direct consequence of the long ballistic flights inherent in the model. This highlights a fundamental mechanistic difference: superdiffusion in the vortex system emerges from a rapidly decorrelating chaotic field with Hamiltonian memory, while in the Lévy walk, it arises from prolonged periods of persistent, memoryless motion.

\begin{figure}[t]
    \centering
    \includegraphics[width=\columnwidth]{../input_files/plots/fig3_vacf.png}
    \caption{Normalized velocity autocorrelation functions (VACF), C$_v$($\Delta$t), for point-vortex tracers (left) and Lévy walks (right) versus lag time. For the point-vortex system, the velocity decorrelates more rapidly as the number of vortices N increases, evidenced by the faster decay of the VACF. The N=40 configuration exhibits a negative lobe, a signature of velocity anti-correlations arising from Hamiltonian memory. In contrast, the superdiffusive Lévy walks (e.g., $\beta=1.2, 1.5$) show a slow, persistent positive correlation, reflecting the long-duration ballistic flights that characterize their motion. The comparison highlights the distinct mechanistic origins of superdiffusion: rapid decorrelation from chaotic interactions in the vortex system versus sustained ballistic motion in the Lévy walk model.}
    \label{fig:vacf}
\end{figure}

\subsection{Lévy-like trapping statistics}

To investigate the origin of the intermittent transport observed in the vortex system, we analyzed the statistics of trapping events, defined as periods where the tracer speed drops below a threshold (the 25th percentile). The probability distributions of the residence times, $P(\tau_{res})$, are shown in Figure \ref{fig:residence_times}.

For all vortex configurations, the distributions exhibit power-law tails, $P(\tau_{res}) \sim \tau_{res}^{-\gamma}$. The key finding is that the tail exponent $\gamma$ decreases as $N$ increases (see Table \ref{tab:residence_exponents}). For $N=40$, we find $\gamma \approx 1.7$. According to continuous-time random walk theory, an exponent $1 < \gamma < 2$ implies a divergent mean residence time, which is the defining feature of Lévy-like trapping. This provides direct mechanistic evidence that the superdiffusion in the dense vortex system is driven by trapping events with a heavy-tailed duration distribution. This result can be connected to the MSD scaling via the relation $\alpha = 2 - (\gamma - 1)$ for $1 < \gamma < 2$. For $N=40$, this predicts an anomalous exponent of $\alpha \approx 2 - (1.7 - 1) = 1.3$, which is in reasonable agreement with the directly measured value of $\alpha \approx 1.52$. In contrast, for $N=10$, $\gamma \approx 2.8 > 2$, and the distribution shows an exponential cutoff, consistent with a finite mean trapping time and near-normal diffusion.

\begin{figure}[t]
    \centering
    \includegraphics[width=\columnwidth]{../input_files/plots/fig4_residence_times.png}
    \caption{Log-log plot of the residence time distributions, $P(\tau_{res})$, for tracers in point-vortex fields with N = 10, 20, and 40. Residence times are identified as periods where the instantaneous tracer speed falls below the 25th percentile. The distributions exhibit power-law tails, $P(\tau_{res}) \sim \tau^{-\gamma}$, with the tail exponent $\gamma$ decreasing as the number of vortices N increases. For the N = 40 configuration, the exponent $\gamma \approx 1.7$ is less than 2, indicating a divergent mean residence time and providing direct evidence for Lévy-like trapping. In contrast, the N = 10 distribution shows an exponential cutoff, consistent with its near-normal diffusion.}
    \label{fig:residence_times}
\end{figure}

\begin{table}[h]
\centering
\caption{Residence time power-law tail exponents.}
\label{tab:residence_exponents}
\begin{tabular}{cc}
\hline
\hline
Number of Vortices ($N$) & Tail Exponent ($\gamma$) \\
\hline
10 & $2.8 \pm 0.4$ \\
20 & $2.2 \pm 0.3$ \\
40 & $1.7 \pm 0.2$ \\
\hline
\end{tabular}
\end{table}

\subsection{Evolution of displacement distributions}

We further examined the one-dimensional displacement probability density functions (PDFs), $P(\Delta x)$, for the most superdiffusive cases: the $N=40$ vortex system and the $\beta=1.5$ Lévy walk. Figure \ref{fig:pdfs} shows the evolution of these PDFs with increasing lag time $\Delta t$.

For both systems, the PDFs are nearly Gaussian at very short lag times. As $\Delta t$ increases, the distributions develop heavy, non-Gaussian tails. At the longest lag times shown, the core of the distributions remains approximately Gaussian, but the tails decay much slower. This convergence towards a non-Gaussian attractor provides further strong evidence that the chaotic vortex flow reproduces the statistical signatures of a Lévy process. Specifically, the tails of the $N=40$ vortex displacement PDF are well-described by a Lévy-stable distribution with stability index $\alpha_{stable} \approx 1.6$, while the $\beta=1.5$ Lévy walk PDF corresponds to $\alpha_{stable} \approx 1.5$. The primary difference observed is that the heavy tails develop more rapidly for the Lévy walk, which is expected since long flights are a fundamental component of the process, whereas in the vortex system they are an emergent property of the chaotic advection.

\begin{figure}[t]
    \centering
    \includegraphics[width=\columnwidth]{../input_files/plots/fig5_displacement_pdfs.png}
    \caption{Marginal displacement probability density functions (PDFs), P($\Delta$x), for the N=40 point-vortex configuration (left) and the $\beta$=1.5 Lévy walk class (right) at four increasing lag times ($\Delta$t). For both systems, the distributions are approximately Gaussian at the shortest lag time but develop progressively heavier, non-Gaussian tails as $\Delta$t increases, shown relative to a Gaussian reference (dashed line). At the largest lag times, both the vortex and Lévy walk PDFs exhibit heavy tails characteristic of a Lévy-stable distribution, providing strong evidence for the effective equivalence of the two systems in terms of their displacement statistics. The Lévy walk PDFs develop heavy tails more rapidly with increasing lag time compared to the vortex tracer PDFs.}
    \label{fig:pdfs}
\end{figure}

\subsection{Ergodicity and its dependence on chaoticity}

While the kinematic analyses suggest a strong correspondence between the two systems, an investigation of their ergodic properties reveals a fundamental divergence. We quantified this using the ergodicity-breaking (EB) parameter, which measures the variance in the time-averaged MSD across different trajectories in an ensemble. A value of $EB=0$ signifies perfect ergodicity.

As shown in Figure \ref{fig:ergodicity}, the point-vortex system becomes progressively \textit{more} ergodic as the number of vortices $N$ increases. The plateau value of the EB parameter decreases from $\approx 0.45$ for $N=5$ to $\approx 0.18$ for $N=40$. This indicates that as the flow becomes more chaotic and superdiffusive, individual tracers sample the phase space more uniformly, leading to greater statistical homogeneity across the ensemble. It is important to note, however, that with only five tracer realizations per configuration, the statistical uncertainty on the EB estimate is large, and the system remains significantly non-ergodic even for $N=40$.

This trend is in stark opposition to the behavior of the Lévy walk model. For Lévy walks, stronger superdiffusion (i.e., smaller $\beta$) leads to stronger ergodicity breaking. The EB parameter increases from $\approx 0.12$ for the normal-diffusive case ($\beta=2.5$) to $\approx 0.62$ for the strongly superdiffusive case ($\beta=1.2$). This occurs because the dynamics become dominated by rare, extremely long flights, causing individual trajectories to have highly variable diffusive properties.

\begin{figure}[t]
    \centering
    \includegraphics[width=\columnwidth]{../input_files/plots/fig6_ergodicity.png}
    \caption{The ergodicity-breaking parameter, EB($\Delta$t), as a function of lag time for point-vortex tracers (left) and Lévy walks (right). The point-vortex system becomes more ergodic (EB decreases) as the number of vortices N increases and the dynamics become more chaotic. In contrast, the Lévy walk system becomes less ergodic (EB increases) as superdiffusion becomes stronger ($\beta$ decreases), highlighting a key mechanistic difference between the two models despite potentially similar mean squared displacement scaling.}
    \label{fig:ergodicity}
\end{figure}

\subsection{The superdiffusion-ergodicity regime map}

The opposing trends in ergodicity are summarized in the regime map of Figure \ref{fig:regime_map}, which plots the anomalous exponent $\alpha$ against the plateau EB parameter for both systems. This plot encapsulates the central finding of our study.

The Lévy walk models trace out a clear, positive correlation: stronger superdiffusion (higher $\alpha$) is intrinsically linked to stronger ergodicity breaking (higher EB). The point-vortex systems, however, occupy a completely different region of this parameter space. For this system, stronger superdiffusion is associated with \textit{weaker} ergodicity breaking.

This result demonstrates that despite the remarkable kinematic similarities—including superdiffusive scaling and heavy-tailed displacement and residence time distributions—the chaotic Hamiltonian flow of the point-vortex system is fundamentally inequivalent to a stochastic Lévy walk. The deterministic constraints and memory effects inherent in the vortex dynamics impose a distinct statistical structure that governs the relationship between transport and ergodicity, precluding a direct mapping to the canonical stochastic model.

\begin{figure}[t]
    \centering
    \includegraphics[width=\columnwidth]{../input_files/plots/fig7_regime_map.png}
    \caption{Regime map of the anomalous exponent $\alpha$ versus the ergodicity-breaking parameter EB for point-vortex systems (blue circles) and Lévy walks (orange squares). The plot illustrates the divergent relationship between superdiffusion and ergodicity breaking in the two models. Lévy walks exhibit a positive correlation, where stronger superdiffusion (higher $\alpha$) corresponds to stronger ergodicity breaking (higher EB). The point-vortex systems occupy a distinct region of the parameter space and do not follow this trend, highlighting a fundamental mechanistic difference between transport driven by chaotic advection and a stochastic walk model.}
    \label{fig:regime_map}
\end{figure}

\end{document}
                